% GENERATED FILE. Edit canonical lessons, then run npm run generate:books.
% canonical-source-hash: fnv1a64:d417242b80223e81
% canonical-lessons: FR-C322-recit, FR-C322-narrer, FR-C322-legende, FR-C322-mythe, FR-C322-morale

\chapter{Le récit}
\label{ch:fr-le-recit}
\hlchaptermodality{322}

\begin{chapteropening}
\textbf{By the end of this chapter:} \emph{I can talk about stories: an account, narrating, a legend, a myth and the moral.}

Complete the last lesson of chapter 322: I can talk about stories: an account, narrating, a legend, a myth and the moral.
\end{chapteropening}

\section[le récit]{le récit — the account, the story}

\label{lesson:FR-C322-recit}



\begin{quote}
Before the new one: say the French for the dragon, then the French for the fable.
\end{quote}

\subsection*{You'll want to know: le récit}
\textbf{le récit} — \textquotedblleft{}the account, the story\textquotedblright{}.

\begin{quote}
Son récit de voyage était très drôle.
\end{quote}

\begin{quote}
\emph{His account of the trip was very funny.}
\end{quote}

\emph{Un récit de voyage} — a travel story.

\begin{cousinweb}
\emph{récit} is built on \emph{réciter} \textquotedblleft{}to recite\textquotedblright{} (EtymDB), a learned borrowing from Latin \emph{recitāre} \textquotedblleft{}to read aloud\textquotedblright{}.
\end{cousinweb}

\begin{grammarlens}[title={What you've built}]
One more word for telling, explaining and describing at B1.
\end{grammarlens}

\subsection*{Your turn}
\begin{itemize}
\raggedright
  \item \emph{Say it:} \emph{le récit}
  \item \emph{Say it:} \emph{le récit}, once more
  \item \emph{Say it:} a sentence of your own with \emph{le récit}
  \item \emph{Recall:} say the French for the dragon, then the French for the fable, then say \emph{le récit} again
\end{itemize}

\begin{tcolorbox}[breakable,colback=teal!4,colframe=teal!35!black,title={Before you move on}]
What does le récit mean? (\textquotedblleft{}The account, the story\textquotedblright{}.) Say it once more.
\end{tcolorbox}
\section[narrer]{narrer — to narrate, to tell}

\label{lesson:FR-C322-narrer}



\begin{quote}
Before the new one: say the French for the fable, then the French for the account.
\end{quote}

\subsection*{You'll want to know: narrer}
\textbf{narrer} — \textquotedblleft{}to narrate, to tell\textquotedblright{}.

\begin{quote}
Le grand-père narrait des histoires de guerre.
\end{quote}

\begin{quote}
\emph{The grandfather used to tell war stories.}
\end{quote}

Literary; in speech, \emph{raconter}.

\begin{cousinweb}
\emph{narrer} is a learned borrowing from Latin \emph{narrāre} \textquotedblleft{}to tell, to relate\textquotedblright{} (EtymDB).
\end{cousinweb}

\begin{grammarlens}[title={What you've built}]
One more word for telling, explaining and describing at B1.
\end{grammarlens}

\subsection*{Your turn}
\begin{itemize}
\raggedright
  \item \emph{Say it:} \emph{narrer}
  \item \emph{Say it:} \emph{narrer}, once more
  \item \emph{Say it:} a sentence of your own with \emph{narrer}
  \item \emph{Recall:} say the French for the fable, then the French for the account, then say \emph{narrer} again
\end{itemize}

\begin{tcolorbox}[breakable,colback=teal!4,colframe=teal!35!black,title={Before you move on}]
What does narrer mean? (\textquotedblleft{}To narrate, to tell\textquotedblright{}.) Say it once more.
\end{tcolorbox}
\section[la légende]{la légende — the legend; the caption}

\label{lesson:FR-C322-legende}



\begin{quote}
Before the new one: say the French for the account, then the French for to narrate.
\end{quote}

\subsection*{You'll want to know: la légende}
\textbf{la légende} — \textquotedblleft{}the legend; the caption\textquotedblright{}.

\begin{quote}
Dans la légende, un dragon vivait dans ce lac.
\end{quote}

\begin{quote}
\emph{In the legend, a dragon lived in this lake.}
\end{quote}

Under a photo or a map, \emph{la légende} is the caption or the key.

\begin{cousinweb}
\emph{légende} is a learned borrowing from Medieval Latin \emph{legenda} \textquotedblleft{}things to be read\textquotedblright{}, from \emph{legere} \textquotedblleft{}to read\textquotedblright{} (Lewis \& Short).
\end{cousinweb}

\begin{grammarlens}[title={What you've built}]
One more word for telling, explaining and describing at B1.
\end{grammarlens}

\subsection*{Your turn}
\begin{itemize}
\raggedright
  \item \emph{Say it:} \emph{la légende}
  \item \emph{Say it:} \emph{la légende}, once more
  \item \emph{Say it:} a sentence of your own with \emph{la légende}
  \item \emph{Recall:} say the French for the account, then the French for to narrate, then say \emph{la légende} again
\end{itemize}

\begin{tcolorbox}[breakable,colback=teal!4,colframe=teal!35!black,title={Before you move on}]
What does la légende mean? (\textquotedblleft{}The legend; the caption\textquotedblright{}.) Say it once more.
\end{tcolorbox}
\section[le mythe]{le mythe — the myth}

\label{lesson:FR-C322-mythe}



\begin{quote}
Before the new one: say the French for to narrate, then the French for the legend.
\end{quote}

\subsection*{You'll want to know: le mythe}
\textbf{le mythe} — \textquotedblleft{}the myth\textquotedblright{}.

\begin{quote}
Les mythes grecs parlent de dieux et de monstres.
\end{quote}

\begin{quote}
\emph{Greek myths talk about gods and monsters.}
\end{quote}

Also a false belief: \emph{c'est un mythe}.

\begin{cousinweb}
\emph{mythe} comes from Late Latin \emph{mythos}, a borrowing from Greek \emph{mûthos} \textquotedblleft{}word, story\textquotedblright{} (EtymDB).
\end{cousinweb}

\begin{grammarlens}[title={What you've built}]
One more word for telling, explaining and describing at B1.
\end{grammarlens}

\subsection*{Your turn}
\begin{itemize}
\raggedright
  \item \emph{Say it:} \emph{le mythe}
  \item \emph{Say it:} \emph{le mythe}, once more
  \item \emph{Say it:} a sentence of your own with \emph{le mythe}
  \item \emph{Recall:} say the French for to narrate, then the French for the legend, then say \emph{le mythe} again
\end{itemize}

\begin{tcolorbox}[breakable,colback=teal!4,colframe=teal!35!black,title={Before you move on}]
What does le mythe mean? (\textquotedblleft{}The myth\textquotedblright{}.) Say it once more.
\end{tcolorbox}
\section[la morale]{la morale — the moral (of a story); morals}

\label{lesson:FR-C322-morale}



\begin{quote}
Before the new one: say the French for the legend, then the French for the myth.
\end{quote}

\subsection*{You'll want to know: la morale}
\textbf{la morale} — \textquotedblleft{}the moral (of a story); morals\textquotedblright{}.

\begin{quote}
La morale de cette histoire est simple.
\end{quote}

\begin{quote}
\emph{The moral of this story is simple.}
\end{quote}

\emph{Faire la morale à quelqu'un} — to lecture someone.

\begin{cousinweb}
\emph{morale} is a learned borrowing from Latin \emph{mōrālis} \textquotedblleft{}of manners\textquotedblright{}, built on \emph{mōs} \textquotedblleft{}custom, habit\textquotedblright{} (Lewis \& Short).
\end{cousinweb}

\begin{grammarlens}[title={What you've built}]
One more word for telling, explaining and describing at B1.
\end{grammarlens}

\subsection*{Your turn}
\begin{itemize}
\raggedright
  \item \emph{Say it:} \emph{la morale}
  \item \emph{Say it:} \emph{la morale}, once more
  \item \emph{Say it:} a sentence of your own with \emph{la morale}
  \item \emph{Recall:} say the French for the legend, then the French for the myth, then say \emph{la morale} again
\end{itemize}

\begin{tcolorbox}[breakable,colback=teal!4,colframe=teal!35!black,title={Before you move on}]
What does la morale mean? (\textquotedblleft{}The moral (of a story); morals\textquotedblright{}.) Say it once more.
\end{tcolorbox}
